\documentclass[12pt,a4paper]{article}

\usepackage[T1]{fontenc}
\usepackage{lmodern}
\usepackage[margin=1in]{geometry}
\usepackage{setspace}
\usepackage{xcolor}
\usepackage[hypertexnames=false]{hyperref}
\usepackage{booktabs}
\usepackage{array}
\usepackage{float}
\usepackage{longtable}
\usepackage{verbatim}
\usepackage{upquote}

\onehalfspacing
\raggedbottom

\definecolor{boxgray}{rgb}{0.93,0.93,0.93}

% A simple shaded definition box
\newcommand{\defbox}[2]{%
  \vspace{0.2cm}
  \noindent\colorbox{boxgray}{\parbox{\dimexpr\linewidth-2\fboxsep}{%
    \textbf{#1:} #2
  }}
  \vspace{0.2cm}
}

\begin{document}

% -------------------------------------------------------------
% Header
% -------------------------------------------------------------
\begin{center}
    {\Large \textbf{CSE312: Database Management Systems Lab}}\\[0.3cm]
    {\large \textbf{Lab 04: PostgreSQL Data Query Language (DQL)}}
\end{center}

\vspace{0.5cm}

% -------------------------------------------------------------
% 1. Objective
% -------------------------------------------------------------
\section*{1. Objective}
The objectives of this experiment are:
\begin{itemize}
    \item To understand the purpose of Data Query Language (DQL).
    \item To retrieve selected columns, expressions, and distinct values with \texttt{SELECT}.
    \item To filter individual rows with a detailed understanding of the \texttt{WHERE} clause.
    \item To sort query results and return a controlled number of rows.
    \item To summarize data using aggregate functions.
    \item To form groups with \texttt{GROUP BY} and filter groups with \texttt{HAVING}.
    \item To combine query results using set operators.
\end{itemize}

% -------------------------------------------------------------
% 2. Introduction
% -------------------------------------------------------------
\section*{2. Introduction to DQL}
\defbox{DQL}{Data Query Language is the part of SQL used to read, filter, calculate, summarize,
and present data without directly changing the stored rows.}

The principal DQL command is \texttt{SELECT}. A simple \texttt{SELECT} reads one table, while an
advanced \texttt{SELECT} may create groups, calculate aggregate values, combine result sets, or
use the result of another query.

\subsection*{2.1 General SELECT Syntax}
\begin{verbatim}
SELECT [DISTINCT] select_list
FROM source
[WHERE row_condition]
[GROUP BY grouping_columns]
[HAVING group_condition]
[ORDER BY sort_expression [ASC | DESC]]
[LIMIT row_count]
[OFFSET rows_to_skip];
\end{verbatim}

\noindent Square brackets in the syntax mean that the enclosed part is optional. They should not
be typed in an actual query.

\noindent\textbf{Complete clause breakdown:}
\begin{itemize}
    \item \texttt{SELECT}: starts the query and specifies the values to display.
    \item \texttt{DISTINCT}: optionally removes duplicate output rows.
    \item \texttt{select\_list}: contains columns, expressions, functions, or subqueries to return.
    \item \texttt{FROM}: identifies the table, view, or subquery that supplies rows.
    \item \texttt{WHERE}: filters individual source rows before grouping.
    \item \texttt{GROUP BY}: collects rows that have equal grouping values.
    \item \texttt{HAVING}: filters complete groups after aggregate values are calculated.
    \item \texttt{ORDER BY}: sorts the final result.
    \item \texttt{ASC}: requests ascending order; it is the default.
    \item \texttt{DESC}: requests descending order.
    \item \texttt{LIMIT}: restricts the maximum number of rows returned.
    \item \texttt{OFFSET}: skips a specified number of rows before returning results.
    \item The semicolon terminates the SQL statement.
\end{itemize}

\subsection*{2.2 Query Clauses at a Glance}
\begingroup
\setstretch{1}
\small
\renewcommand{\arraystretch}{1.25}
\begin{longtable}{>{\bfseries}p{2.5cm} p{4.2cm} p{5.8cm}}
\caption{Important DQL clauses and keywords}
\label{tab:dql_clauses}\\
\toprule
Part & Function & Example \\
\midrule
\endfirsthead
\toprule
Part & Function & Example \\
\midrule
\endhead
\bottomrule
\endfoot
\texttt{SELECT} & Chooses the output columns or expressions. &
\texttt{SELECT first\_name, admission\_year} \\

\texttt{DISTINCT} & Removes duplicate rows from the selected result. &
\texttt{SELECT DISTINCT dept\_id} \\

\texttt{AS} & Assigns a readable alias to a column or source. &
\texttt{COUNT(*) AS total} \\

\texttt{FROM} & Identifies the source of the rows. &
\texttt{FROM student} \\

\texttt{WHERE} & Filters individual rows. &
\texttt{WHERE admission\_year >= 2022} \\

\texttt{GROUP BY} & Forms one group for each distinct grouping value. &
\texttt{GROUP BY dept\_id} \\

\texttt{HAVING} & Filters groups, usually using an aggregate. &
\texttt{HAVING COUNT(*) >= 2} \\

\texttt{ORDER BY} & Sorts the result. &
\texttt{ORDER BY admission\_year DESC} \\

\texttt{LIMIT} & Sets the maximum number of returned rows. &
\texttt{LIMIT 5} \\

\texttt{OFFSET} & Skips rows in the sorted result. &
\texttt{OFFSET 5} \\

\texttt{UNION} & Combines results and removes duplicates. &
\texttt{query\_1 UNION query\_2} \\

\texttt{UNION ALL} & Combines results and preserves duplicates. &
\texttt{query\_1 UNION ALL query\_2} \\

\texttt{INTERSECT} & Returns rows present in both results. &
\texttt{query\_1 INTERSECT query\_2} \\

\texttt{EXCEPT} & Returns rows from the first result that are absent from the second. &
\texttt{query\_1 EXCEPT query\_2} \\
\end{longtable}
\endgroup

\subsection*{2.3 Written Order and Logical Processing Order}
SQL is written in a human-readable order, but its clauses are conceptually processed in a
different order.

\begin{table}[H]
\centering
\renewcommand{\arraystretch}{1.35}
\begin{tabular}{>{\bfseries}p{1.5cm} p{3.2cm} p{7.7cm}}
\toprule
Step & Clause & Conceptual operation \\
\midrule
1 & \texttt{FROM} & Build the source row set. \\
2 & \texttt{WHERE} & Keep source rows whose row condition is true. \\
3 & \texttt{GROUP BY} & Divide the remaining rows into groups. \\
4 & \texttt{HAVING} & Keep groups whose group condition is true. \\
5 & \texttt{SELECT} & Calculate and produce the requested output values. \\
6 & \texttt{DISTINCT} & Remove duplicate output rows when requested. \\
7 & \texttt{ORDER BY} & Sort the output rows. \\
8 & \texttt{OFFSET} & Skip rows from the beginning of the sorted result. \\
9 & \texttt{LIMIT} & Return no more than the requested number of rows. \\
\bottomrule
\end{tabular}
\caption{Conceptual processing order of a SELECT query}
\label{tab:logical_order}
\end{table}

\subsection*{2.4 General Query Safety and Readability Rules}
\begin{itemize}
    \item List only the columns that are required instead of using \texttt{*} in final reports.
    \item Qualify a column with a table alias when more than one source contains that column name.
    \item Use parentheses to make mixed \texttt{AND} and \texttt{OR} conditions unambiguous.
    \item Use \texttt{IS NULL} or \texttt{IS NOT NULL} to test null values.
    \item Use \texttt{ORDER BY} with \texttt{LIMIT} when the identity of the selected rows matters.
    \item Give every subquery in the \texttt{FROM} clause an alias.
    \item Test an inner query separately before placing it inside a larger query.
\end{itemize}

% -------------------------------------------------------------
% 3. Preparing the Database
% -------------------------------------------------------------
\section*{3. Preparing the Database}
This experiment uses the \texttt{university\_lab} database and the four related tables introduced
in Labs 02 and 03.

\begin{verbatim}
psql -U postgres
\c university_lab
\end{verbatim}

If the tables are not available, create them with the following statements.

\begin{verbatim}
CREATE TABLE IF NOT EXISTS department (
    dept_id SERIAL PRIMARY KEY,
    dept_name VARCHAR(100) NOT NULL UNIQUE,
    location VARCHAR(100)
);

CREATE TABLE IF NOT EXISTS student (
    student_id SERIAL PRIMARY KEY,
    first_name VARCHAR(50) NOT NULL,
    last_name VARCHAR(50) NOT NULL,
    email VARCHAR(100) UNIQUE,
    date_of_birth DATE,
    admission_year INTEGER CHECK (admission_year >= 2000),
    is_active BOOLEAN DEFAULT TRUE,
    dept_id INTEGER REFERENCES department(dept_id)
);

CREATE TABLE IF NOT EXISTS course (
    course_id SERIAL PRIMARY KEY,
    course_code VARCHAR(20) NOT NULL UNIQUE,
    course_title VARCHAR(150) NOT NULL,
    credit NUMERIC(3, 1) CHECK (credit > 0),
    dept_id INTEGER NOT NULL REFERENCES department(dept_id)
);

CREATE TABLE IF NOT EXISTS enrollment (
    student_id INTEGER REFERENCES student(student_id),
    course_id INTEGER REFERENCES course(course_id),
    enrollment_date DATE DEFAULT CURRENT_DATE,
    grade VARCHAR(2),
    PRIMARY KEY (student_id, course_id),
    CONSTRAINT grade_check
        CHECK (grade IN ('A+', 'A', 'B+', 'B', 'C+', 'C', 'D', 'F'))
);
\end{verbatim}

\subsection*{3.1 Loading Repeatable Sample Data}
The following data gives the queries in this manual a predictable starting point. It deliberately
includes repeated values, a student without a department, inactive students, and unassigned
grades.

\begin{verbatim}
TRUNCATE TABLE enrollment, student, course, department
RESTART IDENTITY CASCADE;

INSERT INTO department (dept_name, location)
VALUES
    ('Computer Science and Engineering', 'Academic Building 1'),
    ('Electrical and Electronic Engineering', 'Academic Building 2'),
    ('Business Administration', 'Academic Building 3'),
    ('Mathematics', 'Academic Building 4');

INSERT INTO student
    (first_name, last_name, email, date_of_birth,
     admission_year, is_active, dept_id)
VALUES
    ('Amina', 'Rahman', 'amina.rahman@example.com',
     '2003-02-12', 2022, TRUE, 1),
    ('Nafis', 'Ahmed', 'nafis.ahmed@example.com',
     '2002-09-25', 2021, TRUE, 1),
    ('Sadia', 'Islam', 'sadia.islam@example.com',
     '2003-07-08', 2022, TRUE, 2),
    ('Tanvir', 'Hasan', 'tanvir.hasan@example.com',
     '2004-01-19', 2023, TRUE, 1),
    ('Maliha', 'Chowdhury', 'maliha.c@example.com',
     '2002-11-03', 2021, TRUE, 3),
    ('Rafi', 'Karim', 'rafi.karim@example.com',
     '2005-04-15', 2024, TRUE, 1),
    ('Nusrat', 'Jahan', 'nusrat.jahan@example.com',
     '2003-12-21', 2023, FALSE, 2),
    ('Imran', 'Hossain', 'imran.h@example.com',
     '2003-05-10', 2022, TRUE, 3),
    ('Farhan', 'Ali', 'farhan.ali@example.com',
     '2001-08-30', 2020, FALSE, 1),
    ('Tania', 'Sultana', 'tania.s@example.com',
     '2005-06-18', 2024, TRUE, NULL);

INSERT INTO course (course_code, course_title, credit, dept_id)
VALUES
    ('CSE312', 'Database Management Systems Lab', 1.5, 1),
    ('CSE311', 'Database Management Systems', 3.0, 1),
    ('CSE220', 'Data Structures', 3.0, 1),
    ('EEE201', 'Electrical Circuits', 3.0, 2),
    ('EEE202', 'Electronic Devices', 3.0, 2),
    ('BUS101', 'Introduction to Business', 3.0, 3),
    ('BUS201', 'Principles of Marketing', 3.0, 3),
    ('MAT101', 'Discrete Mathematics', 3.0, 4);

INSERT INTO enrollment
    (student_id, course_id, enrollment_date, grade)
VALUES
    (1, 1, '2026-01-15', 'A'),
    (1, 2, '2026-01-15', 'B+'),
    (2, 1, '2026-01-16', 'A+'),
    (2, 3, '2026-01-16', 'A'),
    (3, 4, '2026-01-17', 'B'),
    (3, 5, '2026-01-17', 'B+'),
    (4, 1, '2026-01-18', NULL),
    (4, 2, '2026-01-18', 'A'),
    (5, 6, '2026-01-19', 'A'),
    (5, 7, '2026-01-19', 'B'),
    (6, 1, '2026-01-20', 'B+'),
    (6, 3, '2026-01-20', 'A+'),
    (7, 4, '2026-01-21', 'C+'),
    (7, 5, '2026-01-21', NULL),
    (8, 6, '2026-01-22', 'B+'),
    (8, 7, '2026-01-22', 'A'),
    (9, 2, '2026-01-23', 'C'),
    (9, 8, '2026-01-23', 'B'),
    (10, 8, '2026-01-24', 'A');
\end{verbatim}

% -------------------------------------------------------------
% 4. Basic SELECT
% -------------------------------------------------------------
\section*{4. Basic Data Retrieval with SELECT}

\subsection*{4.1 Selecting All Columns}
\defbox{SELECT}{The \texttt{SELECT} command retrieves rows and calculates the values that will
appear in a query result.}

\noindent\textbf{General syntax:}
\begin{verbatim}
SELECT *
FROM table_name;
\end{verbatim}

\begin{verbatim}
SELECT *
FROM student;
\end{verbatim}

\noindent\textbf{Clause breakdown:}
\begin{itemize}
    \item \texttt{SELECT}: begins the retrieval operation.
    \item \texttt{*}: requests every column from the source.
    \item \texttt{FROM}: introduces the data source.
    \item \texttt{student}: is the table that supplies the rows.
    \item The semicolon ends the statement.
\end{itemize}

\noindent Using \texttt{*} is convenient for exploration. An explicit column list is usually
clearer in reports and application code because it prevents unexpected output changes when the
table structure changes.

\subsection*{4.2 Selecting Particular Columns and Assigning Aliases}
\begin{verbatim}
SELECT
    student_id AS id,
    first_name AS given_name,
    last_name AS family_name
FROM student;
\end{verbatim}

\noindent\textbf{Clause breakdown:}
\begin{itemize}
    \item The comma-separated expressions after \texttt{SELECT} form the output column list.
    \item \texttt{student\_id}, \texttt{first\_name}, and \texttt{last\_name} are source columns.
    \item \texttt{AS}: assigns an output alias without renaming the stored table column.
    \item \texttt{id}, \texttt{given\_name}, and \texttt{family\_name} are headings in the result.
    \item \texttt{FROM student}: reads rows from the \texttt{student} table.
\end{itemize}

\subsection*{4.3 Selecting Expressions}
\begin{verbatim}
SELECT
    student_id,
    first_name || ' ' || last_name AS full_name,
    2026 - admission_year AS years_since_admission
FROM student;
\end{verbatim}

\noindent\textbf{Expression breakdown:}
\begin{itemize}
    \item \texttt{||}: concatenates text values in PostgreSQL.
    \item \texttt{' '}: supplies one space between the first and last names.
    \item \texttt{AS full\_name}: labels the calculated text expression.
    \item \texttt{2026 - admission\_year}: calculates a numeric value for each row.
    \item Expressions affect only the query output; they do not update stored data.
\end{itemize}

\subsection*{4.4 Removing Duplicate Results with DISTINCT}
\noindent\textbf{General syntax:}
\begin{verbatim}
SELECT DISTINCT column_list
FROM table_name;
\end{verbatim}

\begin{verbatim}
SELECT DISTINCT admission_year
FROM student
ORDER BY admission_year;
\end{verbatim}

\noindent\textbf{Clause breakdown:}
\begin{itemize}
    \item \texttt{DISTINCT}: removes duplicate combinations from the complete selected column list.
    \item \texttt{admission\_year}: is the value whose distinct occurrences are displayed.
    \item \texttt{ORDER BY admission\_year}: sorts the unique years in ascending order.
    \item If several columns follow \texttt{DISTINCT}, PostgreSQL compares the complete row of
    selected values, not each column independently.
\end{itemize}

% -------------------------------------------------------------
% 5. WHERE
% -------------------------------------------------------------
\section*{5. Filtering Rows with the WHERE Clause}
\defbox{WHERE}{The \texttt{WHERE} clause evaluates a true-or-false condition for every source row
and keeps only rows for which the condition is true.}

\subsection*{5.1 Basic Comparison Conditions}
\noindent\textbf{General syntax:}
\begin{verbatim}
SELECT column_list
FROM table_name
WHERE condition;
\end{verbatim}

\begin{verbatim}
SELECT student_id, first_name, last_name, admission_year
FROM student
WHERE admission_year >= 2022;
\end{verbatim}

\noindent\textbf{Clause breakdown:}
\begin{itemize}
    \item \texttt{SELECT}: chooses the columns to display.
    \item \texttt{FROM student}: identifies all student rows as the initial source.
    \item \texttt{WHERE}: starts the row-filtering condition.
    \item \texttt{admission\_year >= 2022}: is true for students admitted in 2022 or later.
    \item Rows for which the condition is false or unknown are excluded.
\end{itemize}

\begingroup
\setstretch{1}
\small
\renewcommand{\arraystretch}{1.25}
\begin{longtable}{>{\bfseries}p{2.5cm} p{4.2cm} p{5.8cm}}
\caption{Common operators used in WHERE conditions}
\label{tab:where_operators}\\
\toprule
Operator & Meaning & Example \\
\midrule
\endfirsthead
\toprule
Operator & Meaning & Example \\
\midrule
\endhead
\bottomrule
\endfoot
\texttt{=} & Equal to & \texttt{dept\_id = 1} \\
\texttt{<>} or \texttt{!=} & Not equal to & \texttt{admission\_year <> 2024} \\
\texttt{>} & Greater than & \texttt{credit > 1.5} \\
\texttt{>=} & Greater than or equal to & \texttt{admission\_year >= 2022} \\
\texttt{<} & Less than & \texttt{admission\_year < 2022} \\
\texttt{<=} & Less than or equal to & \texttt{credit <= 3.0} \\
\texttt{AND} & Both connected conditions must be true & \texttt{dept\_id = 1 AND is\_active} \\
\texttt{OR} & At least one connected condition must be true & \texttt{dept\_id = 1 OR dept\_id = 2} \\
\texttt{NOT} & Reverses a condition & \texttt{NOT is\_active} \\
\texttt{BETWEEN} & Tests an inclusive range & \texttt{admission\_year BETWEEN 2021 AND 2023} \\
\texttt{IN} & Tests membership in a value list or query result & \texttt{dept\_id IN (1, 2)} \\
\texttt{LIKE} & Matches a case-sensitive text pattern & \texttt{first\_name LIKE 'A\%'} \\
\texttt{ILIKE} & Matches a case-insensitive text pattern & \texttt{email ILIKE '\%@example.com'} \\
\texttt{IS NULL} & Tests for a null value & \texttt{grade IS NULL} \\
\texttt{IS NOT NULL} & Tests for a non-null value & \texttt{dept\_id IS NOT NULL} \\
\end{longtable}
\endgroup

\subsection*{5.2 Combining Conditions with AND, OR, and NOT}
\begin{verbatim}
SELECT student_id, first_name, admission_year, dept_id
FROM student
WHERE dept_id = 1
  AND (admission_year >= 2022 OR is_active = FALSE);
\end{verbatim}

\noindent\textbf{Condition breakdown:}
\begin{itemize}
    \item \texttt{dept\_id = 1}: requires membership in department 1.
    \item The parenthesized expression accepts either a recent admission year or an inactive
    status.
    \item \texttt{AND}: requires the department condition and the parenthesized condition.
    \item \texttt{OR}: requires at least one condition inside the parentheses.
    \item \texttt{AND} has higher precedence than \texttt{OR}. Parentheses make the intended order
    explicit and reduce mistakes.
\end{itemize}

The following query demonstrates \texttt{NOT}.

\begin{verbatim}
SELECT student_id, first_name, is_active
FROM student
WHERE NOT is_active;
\end{verbatim}

\noindent Because \texttt{is\_active} is a Boolean column, \texttt{NOT is\_active} selects rows
whose value is \texttt{FALSE}.

\subsection*{5.3 Filtering a Range with BETWEEN}
\begin{verbatim}
SELECT student_id, first_name, admission_year
FROM student
WHERE admission_year BETWEEN 2021 AND 2023;
\end{verbatim}

\noindent\textbf{Clause breakdown:}
\begin{itemize}
    \item \texttt{BETWEEN 2021 AND 2023}: includes both boundary values.
    \item The condition is equivalent to
    \texttt{admission\_year >= 2021 AND admission\_year <= 2023}.
    \item \texttt{NOT BETWEEN} may be used to select values outside the inclusive range.
\end{itemize}

\subsection*{5.4 Testing Membership with IN}
\begin{verbatim}
SELECT course_code, course_title, dept_id
FROM course
WHERE dept_id IN (1, 3);
\end{verbatim}

\noindent\textbf{Clause breakdown:}
\begin{itemize}
    \item \texttt{IN}: compares one expression with every value in a list.
    \item \texttt{(1, 3)}: is the list of accepted department identifiers.
    \item The condition is equivalent to \texttt{dept\_id = 1 OR dept\_id = 3}.
    \item \texttt{NOT IN}: requests values that are not members of the list, but null values require
    special care.
\end{itemize}

\subsection*{5.5 Matching Text with LIKE and ILIKE}
\begin{verbatim}
SELECT student_id, first_name, email
FROM student
WHERE first_name ILIKE 'a%';
\end{verbatim}

\noindent\textbf{Pattern breakdown:}
\begin{itemize}
    \item \texttt{ILIKE}: performs case-insensitive pattern matching in PostgreSQL.
    \item \texttt{'a\%'}: matches text that begins with the letter \texttt{a}.
    \item The percent symbol \texttt{\%} matches zero or more characters.
    \item The underscore \texttt{\_} pattern symbol matches exactly one character.
    \item \texttt{LIKE} performs case-sensitive matching under common PostgreSQL collations.
\end{itemize}

\subsection*{5.6 Testing NULL Values}
\begin{verbatim}
SELECT student_id, first_name, dept_id
FROM student
WHERE dept_id IS NULL;
\end{verbatim}

\noindent\textbf{Clause breakdown:}
\begin{itemize}
    \item \texttt{NULL}: represents missing, unknown, or inapplicable information.
    \item \texttt{IS NULL}: tests whether the value is null.
    \item \texttt{IS NOT NULL}: tests whether a known value exists.
    \item \texttt{dept\_id = NULL} is not a valid replacement because ordinary comparison with
    null produces an unknown result.
\end{itemize}

\noindent SQL expresses what condition the rows must satisfy. PostgreSQL decides how to locate and
process those rows efficiently.

% -------------------------------------------------------------
% 6. ORDER BY and LIMIT
% -------------------------------------------------------------
\section*{6. Sorting and Limiting Query Results}

\subsection*{6.1 Sorting with ORDER BY}
\noindent\textbf{General syntax:}
\begin{verbatim}
SELECT column_list
FROM table_name
ORDER BY sort_expression [ASC | DESC], ...;
\end{verbatim}

\begin{verbatim}
SELECT student_id, first_name, last_name, admission_year
FROM student
ORDER BY admission_year DESC, first_name ASC;
\end{verbatim}

\noindent\textbf{Clause breakdown:}
\begin{itemize}
    \item \texttt{ORDER BY}: sorts the final rows.
    \item \texttt{admission\_year DESC}: places later admission years before earlier years.
    \item \texttt{first\_name ASC}: breaks ties by first name in ascending order.
    \item \texttt{ASC}: is optional because ascending order is the default.
    \item Without \texttt{ORDER BY}, PostgreSQL does not guarantee a particular row order.
\end{itemize}

\subsection*{6.2 Returning a Fixed Number of Rows with LIMIT}
\defbox{LIMIT}{The \texttt{LIMIT} clause sets the maximum number of rows that a query may return.}

\noindent\textbf{General syntax:}
\begin{verbatim}
SELECT column_list
FROM table_name
ORDER BY sort_expression
LIMIT row_count;
\end{verbatim}

\begin{verbatim}
SELECT student_id, first_name, admission_year
FROM student
ORDER BY admission_year DESC, student_id
LIMIT 3;
\end{verbatim}

\noindent\textbf{Clause breakdown:}
\begin{itemize}
    \item \texttt{ORDER BY admission\_year DESC}: places the newest admission years first.
    \item \texttt{student\_id}: provides a stable tie-breaker for students with the same year.
    \item \texttt{LIMIT 3}: returns no more than the first three rows of the sorted result.
    \item \texttt{LIMIT} does not mean ``top'' by itself; the meaning of ``first'' is defined by
    \texttt{ORDER BY}.
\end{itemize}

\subsection*{6.3 Skipping Rows with OFFSET}
\begin{verbatim}
SELECT student_id, first_name, admission_year
FROM student
ORDER BY student_id
LIMIT 3 OFFSET 3;
\end{verbatim}

\noindent\textbf{Clause breakdown:}
\begin{itemize}
    \item \texttt{ORDER BY student\_id}: establishes a repeatable order.
    \item \texttt{OFFSET 3}: skips the first three sorted rows.
    \item \texttt{LIMIT 3}: returns at most the next three rows.
    \item This pattern can implement simple page-based output.
    \item Large offsets can become inefficient because PostgreSQL still has to identify and skip
    the earlier rows.
\end{itemize}

% -------------------------------------------------------------
% 7. Aggregate Functions
% -------------------------------------------------------------
\section*{7. Summarizing Data with Aggregate Functions}
\defbox{Aggregate Function}{An aggregate function receives values from several rows and returns
one summary value for the complete set or for each group.}

\subsection*{7.1 Common Aggregate Functions}
\begingroup
\setstretch{1}
\small
\renewcommand{\arraystretch}{1.25}
\begin{longtable}{>{\bfseries}p{2.9cm} p{4.3cm} p{5.3cm}}
\caption{Common aggregate functions}
\label{tab:aggregate_functions}\\
\toprule
Function & Purpose & Example \\
\midrule
\endfirsthead
\toprule
Function & Purpose & Example \\
\midrule
\endhead
\bottomrule
\endfoot
\texttt{COUNT(*)} & Counts rows, including rows containing null values. &
\texttt{COUNT(*) AS student\_count} \\
\texttt{COUNT(column)} & Counts non-null values in one column. &
\texttt{COUNT(grade) AS graded\_count} \\
\texttt{COUNT(DISTINCT column)} & Counts distinct non-null values. &
\texttt{COUNT(DISTINCT dept\_id)} \\
\texttt{SUM(column)} & Adds numeric values. &
\texttt{SUM(credit)} \\
\texttt{AVG(column)} & Calculates the arithmetic mean of non-null numeric values. &
\texttt{AVG(credit)} \\
\texttt{MIN(column)} & Returns the minimum non-null value. &
\texttt{MIN(admission\_year)} \\
\texttt{MAX(column)} & Returns the maximum non-null value. &
\texttt{MAX(admission\_year)} \\
\end{longtable}
\endgroup

\subsection*{7.2 Aggregating the Complete Table}
\begin{verbatim}
SELECT
    COUNT(*) AS total_students,
    COUNT(dept_id) AS students_with_department,
    MIN(admission_year) AS earliest_year,
    MAX(admission_year) AS latest_year
FROM student;
\end{verbatim}

\noindent\textbf{Clause breakdown:}
\begin{itemize}
    \item No \texttt{GROUP BY} appears, so the complete filtered table is treated as one group.
    \item \texttt{COUNT(*)}: counts every student row.
    \item \texttt{COUNT(dept\_id)}: ignores rows whose department identifier is null.
    \item \texttt{MIN} and \texttt{MAX}: find the smallest and largest admission years.
    \item Each \texttt{AS} alias gives a descriptive heading to a calculated result.
\end{itemize}

\subsection*{7.3 Filtering Before Aggregation}
\begin{verbatim}
SELECT COUNT(*) AS active_cse_students
FROM student
WHERE dept_id = 1 AND is_active = TRUE;
\end{verbatim}

\noindent The \texttt{WHERE} clause first keeps active department 1 students. \texttt{COUNT(*)}
then counts only those remaining rows.

% -------------------------------------------------------------
% 8. GROUP BY and HAVING
% -------------------------------------------------------------
\section*{8. Grouping Rows with GROUP BY and HAVING}

\subsection*{8.1 Creating Groups with GROUP BY}
\defbox{GROUP BY}{The \texttt{GROUP BY} clause collects rows with equal grouping values so that
aggregate functions can calculate one result per group.}

\noindent\textbf{General syntax:}
\begin{verbatim}
SELECT grouping_column, aggregate_function(...)
FROM table_name
[WHERE row_condition]
GROUP BY grouping_column
[ORDER BY sort_expression];
\end{verbatim}

\begin{verbatim}
SELECT
    dept_id,
    COUNT(*) AS student_count,
    MIN(admission_year) AS earliest_admission
FROM student
GROUP BY dept_id
ORDER BY dept_id;
\end{verbatim}

\noindent\textbf{Clause breakdown:}
\begin{itemize}
    \item \texttt{FROM student}: supplies the source rows.
    \item \texttt{GROUP BY dept\_id}: creates one group for each distinct department identifier,
    including a separate null group.
    \item \texttt{dept\_id}: may appear in \texttt{SELECT} because it is a grouping column.
    \item \texttt{COUNT(*)}: counts rows inside each department group.
    \item \texttt{MIN(admission\_year)}: calculates one earliest year for each group.
    \item \texttt{ORDER BY dept\_id}: sorts the group result.
\end{itemize}

\noindent In a grouped query, every selected expression should normally be either:
\begin{itemize}
    \item included in the \texttt{GROUP BY} clause, or
    \item calculated by an aggregate function.
\end{itemize}

\subsection*{8.2 Grouping by More Than One Column}
\begin{verbatim}
SELECT dept_id, admission_year, COUNT(*) AS student_count
FROM student
GROUP BY dept_id, admission_year
ORDER BY dept_id, admission_year;
\end{verbatim}

\noindent\textbf{Clause breakdown:}
\begin{itemize}
    \item \texttt{GROUP BY dept\_id, admission\_year}: forms a group for each distinct combination.
    \item Students in the same department but different admission years belong to different groups.
    \item \texttt{COUNT(*)}: reports the size of each combination group.
    \item Both non-aggregate selected columns are listed in \texttt{GROUP BY}.
\end{itemize}

\subsection*{8.3 Filtering Groups with HAVING}
\defbox{HAVING}{The \texttt{HAVING} clause evaluates a condition for each completed group and keeps
only groups for which the condition is true.}

\noindent\textbf{General syntax:}
\begin{verbatim}
SELECT grouping_column, aggregate_function(...)
FROM table_name
[WHERE row_condition]
GROUP BY grouping_column
HAVING group_condition
[ORDER BY sort_expression];
\end{verbatim}

\begin{verbatim}
SELECT course_id, COUNT(*) AS enrollment_count
FROM enrollment
GROUP BY course_id
HAVING COUNT(*) >= 2
ORDER BY enrollment_count DESC, course_id;
\end{verbatim}

\noindent\textbf{Clause breakdown:}
\begin{itemize}
    \item \texttt{GROUP BY course\_id}: creates one group for each course.
    \item \texttt{COUNT(*)}: calculates the number of enrollments in each course group.
    \item \texttt{HAVING COUNT(*) >= 2}: keeps only course groups containing at least two rows.
    \item \texttt{ORDER BY enrollment\_count DESC}: sorts by the aggregate alias from largest to
    smallest.
    \item \texttt{course\_id}: breaks ties in ascending order.
\end{itemize}

\subsection*{8.4 WHERE Compared with HAVING}
\begin{verbatim}
SELECT dept_id, COUNT(*) AS active_student_count
FROM student
WHERE is_active = TRUE
GROUP BY dept_id
HAVING COUNT(*) >= 2
ORDER BY dept_id;
\end{verbatim}

\noindent\textbf{Clause breakdown:}
\begin{itemize}
    \item \texttt{WHERE is\_active = TRUE}: removes inactive rows before groups are formed.
    \item \texttt{GROUP BY dept\_id}: groups only the active rows.
    \item \texttt{HAVING COUNT(*) >= 2}: removes groups containing fewer than two active students.
    \item \texttt{WHERE} answers ``which rows may enter a group?''
    \item \texttt{HAVING} answers ``which completed groups may enter the result?''
\end{itemize}

\begin{table}[H]
\centering
\renewcommand{\arraystretch}{1.35}
\begin{tabular}{>{\bfseries}p{3cm} p{4.5cm} p{4.5cm}}
\toprule
Feature & WHERE & HAVING \\
\midrule
Works on & Individual source rows & Completed groups \\
Logical timing & Before \texttt{GROUP BY} & After \texttt{GROUP BY} \\
Aggregate use & Normally not allowed & Common and expected \\
Typical condition & \texttt{is\_active = TRUE} & \texttt{COUNT(*) >= 2} \\
Performance role & Reduces rows before grouping & Reduces groups after aggregation \\
\bottomrule
\end{tabular}
\caption{Comparison of WHERE and HAVING}
\label{tab:where_having}
\end{table}

% -------------------------------------------------------------
% 9. Expressions
% -------------------------------------------------------------
\section*{9. Conditional and Null-Handling Expressions}

\subsection*{9.1 Producing Categories with CASE}
\noindent\textbf{General syntax:}
\begin{verbatim}
CASE
    WHEN condition_1 THEN result_1
    WHEN condition_2 THEN result_2
    ELSE default_result
END
\end{verbatim}

\noindent\textbf{Example usage:}
\begin{verbatim}
SELECT
    student_id,
    first_name,
    admission_year,
    CASE
        WHEN admission_year >= 2023 THEN 'Recent'
        WHEN admission_year >= 2021 THEN 'Continuing'
        ELSE 'Senior'
    END AS student_category
FROM student
ORDER BY student_id;
\end{verbatim}

\noindent\textbf{Expression breakdown:}
\begin{itemize}
    \item \texttt{CASE}: begins a conditional expression.
    \item Each \texttt{WHEN}: states a condition tested from top to bottom.
    \item \texttt{THEN}: supplies the value for the first true condition.
    \item \texttt{ELSE}: supplies a value when no \texttt{WHEN} condition is true.
    \item \texttt{END}: closes the expression.
\end{itemize}

\subsection*{9.2 Replacing NULL for Display with COALESCE}
\begin{verbatim}
SELECT
    student_id,
    first_name,
    COALESCE(email, 'No email recorded') AS contact_email
FROM student
ORDER BY student_id;
\end{verbatim}

\noindent\textbf{Expression breakdown:}
\begin{itemize}
    \item \texttt{COALESCE(value\_1, value\_2, ...)}: returns the first non-null argument.
    \item \texttt{email}: is used when the student record contains an email address.
    \item \texttt{'No email recorded'}: is returned when the email value is null.
    \item \texttt{COALESCE} changes the displayed result, not the stored value.
\end{itemize}

% -------------------------------------------------------------
% 10. Set operations
% -------------------------------------------------------------
\section*{10. Combining Query Results with Set Operators}
Set operators combine complete query results vertically. The participating queries must return the
same number of columns, and corresponding columns must have compatible data types.

\subsection*{10.1 UNION and UNION ALL}
\begin{verbatim}
SELECT first_name, last_name
FROM student
WHERE dept_id = 1

UNION

SELECT first_name, last_name
FROM student
WHERE admission_year = 2024
ORDER BY first_name, last_name;
\end{verbatim}

\noindent\textbf{Clause breakdown:}
\begin{itemize}
    \item The first \texttt{SELECT}: returns department 1 student names.
    \item The second \texttt{SELECT}: returns students admitted in 2024.
    \item \texttt{UNION}: combines both results and removes duplicate rows.
    \item \texttt{UNION ALL}: may replace \texttt{UNION} when duplicates should be preserved.
    \item The final \texttt{ORDER BY}: sorts the complete combined result.
    \item Output column names come from the first query.
\end{itemize}

\subsection*{10.2 INTERSECT}
\begin{verbatim}
SELECT student_id
FROM student
WHERE dept_id = 1

INTERSECT

SELECT student_id
FROM enrollment
WHERE course_id = 1;
\end{verbatim}

\noindent\textbf{Clause breakdown:}
\begin{itemize}
    \item The first query returns identifiers of department 1 students.
    \item The second query returns identifiers of students enrolled in course 1.
    \item \texttt{INTERSECT}: keeps identifiers that occur in both query results.
\end{itemize}

\subsection*{10.3 EXCEPT}
\begin{verbatim}
SELECT student_id
FROM student

EXCEPT

SELECT student_id
FROM enrollment;
\end{verbatim}

\noindent\textbf{Clause breakdown:}
\begin{itemize}
    \item The first query returns every student identifier.
    \item The second query returns identifiers that appear in enrollment records.
    \item \texttt{EXCEPT}: keeps students present in the first result but absent from the second.
    \item Reversing the query order changes the meaning of \texttt{EXCEPT}.
\end{itemize}

% -------------------------------------------------------------
% 11. Complete query
% -------------------------------------------------------------
\section*{11. Reading a Complete DQL Query}
The following query combines row filtering, grouping, group filtering, sorting, and limiting.

\begin{verbatim}
SELECT
    dept_id,
    admission_year,
    COUNT(*) AS active_students
FROM student
WHERE is_active = TRUE
  AND dept_id IS NOT NULL
GROUP BY dept_id, admission_year
HAVING COUNT(*) >= 2
ORDER BY active_students DESC, dept_id, admission_year
LIMIT 3;
\end{verbatim}

\noindent\textbf{Clause-by-clause interpretation:}
\begin{enumerate}
    \item \texttt{FROM student}: supplies the source rows.
    \item \texttt{WHERE is\_active = TRUE}: keeps only active students.
    \item \texttt{dept\_id IS NOT NULL}: removes students without an assigned department.
    \item \texttt{GROUP BY dept\_id, admission\_year}: creates one group for each department and
    admission-year combination.
    \item \texttt{COUNT(*) AS active\_students}: counts the rows inside each group.
    \item \texttt{HAVING COUNT(*) >= 2}: keeps only groups with at least two active students.
    \item \texttt{ORDER BY}: ranks larger groups first, then applies stable tie-breakers.
    \item \texttt{LIMIT 3}: returns at most the first three ranked groups.
\end{enumerate}

% -------------------------------------------------------------
% 12. Observations
% -------------------------------------------------------------
\section*{12. Observations}
\begin{itemize}
    \item \texttt{SELECT} retrieved stored columns and calculated expressions without changing
    table data.
    \item \texttt{WHERE} filtered individual rows before grouping.
    \item \texttt{ORDER BY} established a defined result order.
    \item \texttt{LIMIT} and \texttt{OFFSET} selected a portion of a sorted result.
    \item Aggregate functions summarized complete tables and groups.
    \item \texttt{GROUP BY} created one result row for each distinct grouping value.
    \item \texttt{HAVING} filtered groups after aggregate values were available.
    \item Set operators combined compatible query results.
    \item Subqueries supplied values, sets, calculated columns, and table-like sources to outer
    queries.
\end{itemize}

% -------------------------------------------------------------
% 13. Lab Tasks
% -------------------------------------------------------------
\section*{13. Lab Task}

\subsection*{Task A: Basic Retrieval and WHERE Practice}
Using the \texttt{university\_lab} database, write and execute queries to:
\begin{enumerate}
    \item Display student identifiers, full names, and admission years.
    \item Display all distinct admission years in descending order.
    \item Find active students admitted from 2021 through 2023.
    \item Find students whose first names begin with \texttt{M}, \texttt{N}, or \texttt{R}.
    \item Find students who do not have a department.
    \item Find courses belonging to departments 1, 2, or 4.
    \item Add a short clause breakdown below each query.
\end{enumerate}

\subsection*{Task B: Sorting, LIMIT, and OFFSET}
\begin{enumerate}
    \item Display the five most recently admitted students. Use a stable tie-breaker.
    \item Display the second page of students when each page contains three rows.
    \item Display the three courses with the highest credit values, sorted by course code when
    credits are equal.
    \item Explain why each query requires \texttt{ORDER BY} before \texttt{LIMIT}.
\end{enumerate}

\subsection*{Task C: Aggregate, GROUP BY, and HAVING}
\begin{enumerate}
    \item Count all students and count only students with assigned departments.
    \item Show the number of students in each department.
    \item Show the number of students for each department and admission year combination.
    \item Show only departments containing at least two active students.
    \item Show courses with at least three enrollments.
    \item Show each course's total enrollment count and number of assigned grades.
    \item Explain which conditions belong in \texttt{WHERE} and which belong in \texttt{HAVING}.
\end{enumerate}

\subsection*{Task D: Online Shop DQL Problem}
Assume an online shop database contains the tables \texttt{customer}, \texttt{product},
\texttt{orders}, and \texttt{order\_item}. Write PostgreSQL queries to:
\begin{itemize}
    \item display products whose price is between two chosen values;
    \item display the five products with the lowest stock;
    \item calculate the number and total value of items in each order;
    \item display only orders whose calculated value exceeds a chosen amount;
    \item display customers who have never placed an order;
\end{itemize}

% -------------------------------------------------------------
% 14. Conclusion
% -------------------------------------------------------------
\section*{14. Conclusion}
This experiment introduced PostgreSQL Data Query Language through progressively structured
\texttt{SELECT} statements. Students selected and calculated output values, filtered rows, sorted
and limited results, summarized data, formed and filtered groups, and combined result sets.

\end{document}
